\chapter{结合代数、Clifford 代数与旋量结构}
\label{chap:algebra-number-systems}

上一章把对称性组织成群与表示。本章转向一个不同但紧密相连的问题：当向量空间还带有乘法、二次型或共轭结构时，怎样把这些数据统一编码为代数？复数、四元数和八元数说明“扩张数系”会逐步牺牲交换律与结合律；Clifford 代数则从二次型出发，通过 泛性质 构造一个结合代数，并由其可逆元素内部地产生正交群及其二重覆盖。这样，Pauli 矩阵、旋量与 gamma 矩阵都成为一般 Clifford 表示的结果，而不再是先验写下的特殊矩阵。

\section{代数结构、共轭与 composition algebra}
\label{sec:algebra-composition}
设 $\mathbb F=\RR$ 或 $\CC$。一个含单位的结合 $\mathbb F$-代数 $A$ 同时是 $\mathbb F$-向量空间和含单位结合环，并满足
\[
\lambda(ab)=(\lambda a)b=a(\lambda b),
\qquad
\lambda\in\mathbb F.
\]
其中心为
\[
Z(A)=\{z\in A:za=az,\ \forall a\in A\}.
\]
若 $Z(A)=\mathbb F1$，则称代数在标量意义下是中心的。

很多物理与几何代数还带有 对合，即映射 $a\mapsto a^*$ 满足
\[
(a+b)^*=a^*+b^*,
\qquad
(ab)^*=b^*a^*,
\qquad
(a^*)^*=a.
\]
复代数通常还要求 $(\lambda a)^*=\overline\lambda a^*$。矩阵代数中的 厄米 共轭、复数共轭以及四元数共轭都是这一结构的例子。

若代数带有二次型或范数 $N:A\to\RR_{\ge0}$，且满足乘法相容性
\begin{equation}
N(ab)=N(a)N(b),
\label{eq:composition-norm}
\end{equation}
则称其具有 composition 性。若 $N(a)>0$ 对所有 $a\neq0$ 成立，并且 $a^{-1}=a^*/N(a)$，则非零元素自动可逆。有限维实数情形中，正定 multiplicative 范数 把可能性强烈限制为 $\RR,\CC,\mathbb H,\mathbb O$ 四类；它们依次展示“交换结合、交换结合、非交换结合、非交换非结合但交替”的层级。
\label{sec:quaternions}

四元数代数写成
\[
\mathbb H
=\{a+b\mathbf i+c\mathbf j+d\mathbf k:
a,b,c,d\in\RR\},
\]
其中
\[
\mathbf i^2=\mathbf j^2=\mathbf k^2=-1,
\qquad
\mathbf i\mathbf j=\mathbf k,
\quad
\mathbf j\mathbf k=\mathbf i,
\quad
\mathbf k\mathbf i=\mathbf j,
\]
交换两个不同虚基元会改变符号。共轭定义为
\[
\overline q=a-b\mathbf i-c\mathbf j-d\mathbf k,
\]
于是
\[
q\overline q=\overline q q
=a^2+b^2+c^2+d^2
\equiv N(q).
\]
直接利用共轭的反自同构性质
\[
\overline{pq}=\overline q\,\overline p
\]
可得
\[
N(pq)
=pq\overline q\,\overline p
=pN(q)\overline p
=N(q)N(p),
\]
所以四元数满足\eqnrefs{eq:composition-norm}。

把纯虚四元数与 $\RR^3$ 识别，若 $\mathbf a,\mathbf b\in\operatorname{Im}\mathbb H$，则
\begin{equation}
\mathbf a\mathbf b
=-\mathbf a\cdot\mathbf b
+\mathbf a\times\mathbf b.
\label{eq:quaternion-dot-cross}
\end{equation}
单位四元数
\[
q=\cos\frac\theta2+\mathbf n\sin\frac\theta2,
\qquad
|\mathbf n|=1,
\]
通过共轭作用
\[
\mathbf v\longmapsto q\mathbf vq^{-1}
\]
给出绕 $\mathbf n$ 的三维转动。半角不是额外规定，而是因为向量作用同时包含左乘和右乘。这个构造稍后会被包含在一般 $\Spin(3)\to SO(3)$ 中。

\begin{derivation}[四元数共轭作用给出 Rodrigues 公式]
记
\[
c=\cos\frac\theta2,
\qquad
s=\sin\frac\theta2,
\qquad
q=c+s\mathbf n,
\qquad
q^{-1}=c-s\mathbf n.
\]
展开得
\[
q\mathbf vq^{-1}
=c^2\mathbf v
+cs(\mathbf n\mathbf v-\mathbf v\mathbf n)
-s^2\mathbf n\mathbf v\mathbf n.
\]
由\eqnrefs{eq:quaternion-dot-cross}，
\[
\mathbf n\mathbf v-\mathbf v\mathbf n
=2\mathbf n\times\mathbf v.
\]
又由 $|\mathbf n|=1$ 与向量三重积公式，
\[
\mathbf n\mathbf v\mathbf n
=
\mathbf v-2\mathbf n(\mathbf n\cdot\mathbf v).
\]
因此
\[
\begin{aligned}
q\mathbf vq^{-1}
&=(c^2-s^2)\mathbf v
+2cs\,\mathbf n\times\mathbf v
+2s^2\mathbf n(\mathbf n\cdot\mathbf v)\\
&=
\mathbf v\cos\theta
+(\mathbf n\times\mathbf v)\sin\theta
+\mathbf n(\mathbf n\cdot\mathbf v)(1-\cos\theta),
\end{aligned}
\]
这就是 Rodrigues 公式。
\end{derivation}
Cayley--Dickson doubling 从带 对合 的代数 $A$ 构造 $A\oplus A$。一种常用约定为
\[
(a,b)(c,d)
=
(ac-d^*b,\ da+bc^*),
\qquad
(a,b)^*=(a^*,-b).
\]
从 $\RR$ 出发依次得到 $\CC$、$\mathbb H$、$\mathbb O$。但每次 doubling 并不保持全部性质：八元数 $\mathbb O$ 已不满足一般结合律，只满足交替律
\[
(xx)y=x(xy),
\qquad
x(yy)=(xy)y.
\]
因此不能把八元数当成普通矩阵代数处理，更不能把任意八元数恒等式无条件搬到结合代数中。

八元数的自同构群是紧致例外 Lie 群 $G_2$。固定一个单位虚八元数 $u$ 后，保持 $u$ 不变的稳定子同构于 $SU(3)$。这给出
\[
SU(3)\subset G_2,
\]
但其逻辑是“$SU(3)$ 作为八元数自同构群中的稳定子”，而不是“Gell--Mann 矩阵是八元数乘法表的矩阵化”。$SU(3)$ 的一般表示论仍属于上一章的 Lie 群框架。

\section{Clifford 代数：由二次型生成的 universal algebra}
\label{sec:clifford-universal}
设 $V$ 是有限维 $\mathbb F$-向量空间，$Q:V\to\mathbb F$ 是二次型。Clifford 代数 $\Cl(V,Q)$ 是含单位结合代数，并带线性映射
\[
i:V\to\Cl(V,Q)
\]
满足
\begin{equation}
i(v)^2=Q(v)1,
\qquad v\in V,
\label{eq:clifford-defining-relation}
\end{equation}
且具有 泛性质：若任意含单位结合代数 $A$ 与线性映射 $f:V\to A$ 也满足
\[
f(v)^2=Q(v)1_A,
\]
则存在唯一代数同态
\[
\widetilde f:\Cl(V,Q)\to A
\]
使 $\widetilde f\circ i=f$。因此 Clifford 代数不是“选一组 gamma 矩阵”定义出来的；矩阵只是在 泛 代数 上选择一个表示。

由二次型极化得到对称双线性型
\[
B(u,v)=\frac12\bigl(Q(u+v)-Q(u)-Q(v)\bigr).
\]
\eqnrefs{eq:clifford-defining-relation} 等价于
\begin{equation}
uv+vu=2B(u,v)1.
\label{eq:clifford-anticommutator}
\end{equation}
若 $\{e_a\}_{a=1}^n$ 为 $B$-正交基，则有
\[
e_ae_b+e_be_a=2B_{ab}1.
\]
每个基元最多出现一次，因而线性空间维数为
\begin{equation}
\dim\Cl(V,Q)=2^n.
\label{eq:clifford-dimension}
\end{equation}

Clifford 代数具有自然 $\mathbb Z_2$-分次
\[
\Cl(V,Q)
=
\Cl^0(V,Q)
\oplus
\Cl^1(V,Q),
\]
其中偶数个向量的乘积生成 偶 subalgebra $\Cl^0$。与外代数不同，Clifford 乘积同时包含对称的度量收缩与反对称的外积；对 $u,v\in V$，
\[
uv=B(u,v)+u\wedge v.
\]
这正是 Clifford 代数能够同时编码 度量 与 定向 multivector 的原因。

\begin{derivation}[由二次型关系推出 Clifford 反对易关系与维数上界]
对任意 $u,v\in V$，由\eqnrefs{eq:clifford-defining-relation} 有
\[
(u+v)^2=Q(u+v)1.
\]
左边展开为
\[
u^2+uv+vu+v^2.
\]
再代入 $u^2=Q(u)1$、$v^2=Q(v)1$，得到
\[
uv+vu
=\bigl(Q(u+v)-Q(u)-Q(v)\bigr)1
=2B(u,v)1,
\]
即\eqnrefs{eq:clifford-anticommutator}。

若 $\{e_a\}$ 是正交基，则对 $a\neq b$，
\[
e_ae_b=-e_be_a,
\]
而
\[
e_a^2=B(e_a,e_a)1
\]
可化成标量。因此任意单项式都可以重排成
\[
e_{a_1}e_{a_2}\cdots e_{a_k},
\qquad
a_1<a_2<\cdots<a_k.
\]
这类有序单项式共有
\[
\sum_{k=0}^n\binom nk=2^n
\]
个，所以 $\dim\Cl(V,Q)\le2^n$。另一方面，可在外代数 $\Lambda V$ 上构造忠实的左 Clifford 作用
\[
c(v)=\varepsilon(v)+\iota_v,
\]
其中 $\varepsilon(v)$ 为外乘、$\iota_v$ 为由 $B$ 定义的收缩，且
\[
c(u)c(v)+c(v)c(u)=2B(u,v)I.
\]
由此有序单项式线性独立，故等号成立，得到\eqnrefs{eq:clifford-dimension}。
\end{derivation}
仅由“所有有序单项式都能化简出来”只能得到维数上界；要真正理解 Clifford 代数的线性结构，还需要说明这些单项式彼此独立。最自然的做法是给张量代数按次数加过滤
\[
F^0\subset F^1\subset F^2\subset\cdots,
\qquad
F^r\Cl(V,Q)
=\operatorname{span}\{v_1\cdots v_k:k\le r\}.
\]
定义关系
\[
uv+vu=2B(u,v)
\]
把二次项换成低两阶的标量，因此在 伴随 分次 代数 中低阶项消失，只留下
\[
uv+vu=0.
\]
于是存在自然满射
\[
\Lambda V\longrightarrow
\operatorname{gr}\Cl(V,Q).
\]
反过来利用有序单项式计数可知两边各次数的维数完全相同，从而得到 PBW 型同构
\begin{equation}
\operatorname{gr}\Cl(V,Q)
\cong \Lambda V.
\label{eq:clifford-pbw-graded}
\end{equation}
这里的同构是分次向量空间意义下的；它绝不意味着 Clifford 乘法等于外积。真正发生的是：Clifford 乘法的最高次数部分就是外积，而较低次数部分记录由度量产生的收缩。

选取基 $e_1,\ldots,e_n$ 后，PBW 结论说明
\begin{equation}
1,
\quad e_{i_1},
\quad e_{i_1}e_{i_2},
\quad\ldots,\quad
e_{i_1}\cdots e_{i_k},
\qquad i_1<\cdots<i_k,
\label{eq:clifford-pbw-basis}
\end{equation}
确实构成一组基，而不只是生成集。特别地
\[
\dim\Cl(V,Q)=\sum_{k=0}^n\binom nk=2^n.
\]

这个结构可以在外代数上具体实现。令 $\varepsilon(v)$ 表示外乘算子
\[
\varepsilon(v)\omega=v\wedge\omega,
\]
并用双线性型 $B$ 定义收缩 $\iota_v$，使得在一次向量上
\[
\iota_v w=B(v,w).
\]
它作为次数 $-1$ 的反导子延拓到整个 $\Lambda V$。于是
\[
\varepsilon(u)\varepsilon(v)+\varepsilon(v)\varepsilon(u)=0,
\]
\[
\iota_u\iota_v+\iota_v\iota_u=0,
\]
以及最关键的 正则 anticommutation 关系
\begin{equation}
\iota_u\varepsilon(v)+\varepsilon(v)\iota_u
=B(u,v)I.
\label{eq:clifford-wedge-contraction-car}
\end{equation}
因此 Chevalley 算子
\begin{equation*}c(v)=\varepsilon(v)+\iota_v\end{equation*}
满足
\[
c(u)c(v)+c(v)c(u)=2B(u,v)I,
\]
由 泛性质 唯一延拓为 Clifford 表示。

\begin{derivation}[从过滤到 PBW，再到外积加收缩]
把张量代数 $T(V)$ 的次数过滤传到商代数
\[
\Cl(V,Q)
=T(V)/\langle v\otimes v-Q(v)1\rangle.
\]
定义 $F^r\Cl=\operatorname{span}\{v_1\cdots v_k:k\le r\}$。关系 $v\otimes v-Q(v)1$ 中，二次项 $v\otimes v$ 与零次项 $Q(v)1$ 差两个过滤阶；进入 伴随 分次
\[
\operatorname{gr}\Cl=\bigoplus_{r}\frac{F^r}{F^{r-1}}
\]
后，零次项落在 $F^0$ 而消失，只剩二次项的像
\[
[v]^2=0.
\]

极化出反对易关系。把 $[u+v]^2=0$ 展开：
\begin{equation*}
  0=[u+v]^2
  =[u]^2+[u][v]+[v][u]+[v]^2
  =[u][v]+[v][u],
\end{equation*}
其中用了 $[u]^2=[v]^2=0$。故在 $\operatorname{gr}\Cl$ 中
\[
[u][v]+[v][u]=0.
\]

外代数 泛性质 给出满射。外代数 $\Lambda V$ 正是「由 $V$ 生成、且元素反对易」的最一般代数。既然 $\operatorname{gr}\Cl$ 中向量乘积反对易，就存在唯一代数同态
\[
\Lambda V\twoheadrightarrow
\operatorname{gr}\Cl(V,Q),
\]
它满射，因为 $\operatorname{gr}\Cl$ 由各个 $[v]$ 的乘积生成，而每个反对易乘积都是某个外积的像。

维数守恒把满射升级为同构。满射给出上界
\[
\dim\operatorname{gr}\Cl
\le\dim\Lambda V=2^n.
\]
另一方面，有限维过滤空间的 伴随 分次 保持维数：
\[
\dim\operatorname{gr}\Cl
=\sum_r\dim\Bigl(F^r\big/F^{r-1}\Bigr)
=\dim\Cl=2^n,
\]
其中 $\dim\Cl=2^n$ 已由\eqnrefs{eq:clifford-dimension} 得到。于是 $\dim\operatorname{gr}\Cl=2^n=\dim\Lambda V$，同维有限空间之间的满射必为同构，得到\eqnrefs{eq:clifford-pbw-graded}。据此\eqnrefs{eq:clifford-pbw-basis} 中的有序单项式不只是生成集，而是真正线性独立的一组基。

验证 正则反交换关系。对可分解 $k$-形式
\[
\omega=w_1\wedge\cdots\wedge w_k,
\]
收缩 $\iota_u$（次数 $-1$ 的反导子，满足 $\iota_u(v)=B(u,v)$）写为
\[
\iota_u\omega
=\sum_{j=1}^k(-1)^{j-1}B(u,w_j)
 w_1\wedge\cdots\widehat{w_j}\cdots\wedge w_k.
\]
把 $u$ 作用在 $v\wedge\omega=v\wedge w_1\wedge\cdots\wedge w_k$ 上，按上述公式逐项收缩（$v$ 记作第 $0$ 个因子）：
\begin{equation*}
 \iota_u(v\wedge\omega)
 =B(u,v)\,\omega
 +\sum_{j=1}^{k}(-1)^{j}B(u,w_j)\,v\wedge w_1\wedge\cdots\widehat{w_j}\cdots\wedge w_k.
\end{equation*}
第二项恰为 $-v\wedge\iota_u\omega$（把 $v$ 提出并把 $(-1)^j$ 写成 $-(-1)^{j-1}$），故
\[
\iota_u(v\wedge\omega)+v\wedge\iota_u\omega
=B(u,v)\omega.
\]
这正是\eqnrefs{eq:clifford-wedge-contraction-car}。因此 Chevalley 算子 $c(v)=\varepsilon(v)+\iota_v$ 确实满足 Clifford 关系；从向量空间层面看，Clifford 乘法等于
\[
\text{外乘} + \text{度量收缩},
\]
而 PBW 同构保证这种「外积加收缩」的描述没有丢失任何线性信息。
\end{derivation}

在复数域、$B$ 非退化时，PBW 结构还与完整矩阵分类兼容：
\[
\Cl_{2m}(\CC)\cong M_{2^m}(\CC),
\qquad
\Cl_{2m+1}(\CC)
\cong M_{2^m}(\CC)\oplus M_{2^m}(\CC).
\]
偶数维代数为 中心 简单 代数；奇数维则由归一化体积元的两个本征值分裂成两个 简单 块。这个分类说明“gamma 矩阵尺寸”为何由维数强制决定，而不是一种任意表示习惯。

Clifford 代数上有三个常用 对合的 operations。级 对合 $\alpha$ 在向量上满足 $\alpha(v)=-v$，并按乘法延拓；reversion $\tau$ 反转乘积次序：
\[
\tau(v_1v_2\cdots v_k)
=v_k\cdots v_2v_1;
\]
Clifford 共轭 定义为
\[
\overline x=\alpha(\tau(x)).
\]
这些操作在定义 Pin/Spin 群、双线性旋量配对和实结构时不可互换。不同教材若把三者都写成“共轭”，往往会造成符号混乱，因此以后明确写出所用操作。
在复数域上，非退化二次型的符号差异可以通过基变换吸收，因此复 Clifford 代数只依赖维数。作为非分次代数，有
\[
\Cl_{2m}(\CC)
\cong M_{2^m}(\CC),
\]
\[
\Cl_{2m+1}(\CC)
\cong M_{2^m}(\CC)\oplus M_{2^m}(\CC).
\]
偶数维情形因此有维数 $2^m$ 的不可约复旋量模。奇数维情形有两个不等价的不可约模，它们由最高次数体积元的本征值区分。

实 Clifford 代数还依赖 符号差 $(p,q)$，并表现出 模-8 periodicity。后续若需要实旋量、实结构或不同 符号差 之间的比较，就必须保留这一区别；仅研究复表示时则不应把实分类的 符号差 依赖混入复 Clifford 结论。

为了把这句话变成可计算的结构，约定
\[
\Cl_{p,q}:=\Cl(\RR^{p+q},Q),
\qquad
Q(x)=x_1^2+\cdots+x_p^2-x_{p+1}^2-\cdots-x_{p+q}^2.
\]
因此前 $p$ 个正交基向量平方为 $+1$，后 $q$ 个平方为 $-1$。设
\[
n=p+q,
\qquad
m=\left\lfloor\frac n2\right\rfloor,
\qquad
r=p-q\pmod 8.
\]
则实 Clifford 代数的非分次分类为
\begin{equation}
\begin{array}{c|c}
r & \Cl_{p,q} \\
\hline
0 & M_{2^m}(\RR)\\
1 & M_{2^m}(\RR)\oplus M_{2^m}(\RR)\\
2 & M_{2^m}(\RR)\\
3 & M_{2^m}(\CC)\\
4 & M_{2^{m-1}}(\mathbb H)\\
5 & M_{2^{m-1}}(\mathbb H)\oplus M_{2^{m-1}}(\mathbb H)\\
6 & M_{2^{m-1}}(\mathbb H)\\
7 & M_{2^m}(\CC)
\end{array}
\label{eq:real-clifford-classification}
\end{equation}
其中只在相应行确实出现时读取矩阵阶数；例如
\[
\Cl_{1,0}\cong\RR\oplus\RR,
\qquad
\Cl_{0,1}\cong\CC,
\qquad
\Cl_{0,2}\cong\mathbb H,
\]
\[
\Cl_{2,0}\cong M_2(\RR),
\qquad
\Cl_{3,0}\cong M_2(\CC),
\qquad
\Cl_{1,3}\cong M_2(\mathbb H).
\]
这张表不仅给出代数维数，还决定不可约实模的标量交换子究竟是 $\RR$、$\CC$ 还是 $\mathbb H$。因此所谓“实型、复型、四元数型”旋量，本质上来自不可约模交换子代数，而不是对 gamma 矩阵额外附加的经验条件。

分类的两个关键递推是
\begin{equation}
\Cl_{p+1,q+1}
\cong M_2\!\left(\Cl_{p,q}\right),
\label{eq:clifford-pq-recursion}
\end{equation}
以及八周期
\begin{equation}
\Cl_{p+8,q}
\cong M_{16}\!\left(\Cl_{p,q}\right),
\qquad
\Cl_{p,q+8}
\cong M_{16}\!\left(\Cl_{p,q}\right).
\label{eq:clifford-bott-eight-periodicity}
\end{equation}
后者正是实 Clifford 理论中的 Bott 模-8 周期在代数层面的有限维影子。与之相对，复 Clifford 代数只有 模-2 周期：
\[
\Cl_{n+2}(\CC)
\cong M_2\!\left(\Cl_{n}(\CC)\right)
\]
在 简单/半单 块 的意义下递推。

\begin{derivation}[从 $\Cl_{1,1}$ 到矩阵递推与八周期]
先看最小的混合 符号差。取
\[
E_+=
\begin{pmatrix}
0&1\\
1&0
\end{pmatrix},
\qquad
E_-=
\begin{pmatrix}
0&1\\
-1&0
\end{pmatrix}.
\]
则
\[
E_+^2=I,
\qquad
E_-^2=-I,
\qquad
E_+E_-+E_-E_+=0.
\]
因此 泛性质 给出
\[
\Cl_{1,1}\longrightarrow M_2(\RR).
\]
两边实维数同为 $4$，而 $I,E_+,E_-,E_+E_-$ 线性独立，故
\[
\Cl_{1,1}\cong M_2(\RR).
\]

对一般 $\Cl_{p,q}$，把原二次空间的每个向量 $v$ 映为
\[
v\longmapsto
\begin{pmatrix}
v&0\\
0&-v
\end{pmatrix},
\]
再把新增的正、负两个生成元分别映成 $E_+$、$E_-$。这些矩阵两两满足正确的 Clifford 反对易关系，于是 泛性质 产生
\[
\Cl_{p+1,q+1}
\longrightarrow M_2(\Cl_{p,q}).
\]
由生成性与维数
\[
2^{p+q+2}=4\cdot2^{p+q}
\]
可知它是同构，得到\eqnrefs{eq:clifford-pq-recursion}。

再配合低维种子
\[
\Cl_{0,1}\cong\CC,
\qquad
\Cl_{0,2}\cong\mathbb H,
\qquad
\Cl_{1,0}\cong\RR\oplus\RR,
\]
以及 张量-积 的分次版本反复递推，就得到\eqnrefs{eq:real-clifford-classification}。走完八步后可除代数类型回到原点，而矩阵阶数累积出因子 $16$，于是得到\eqnrefs{eq:clifford-bott-eight-periodicity}。这里真正周期的是 Morita 类型；矩阵阶数负责补偿总维数的增长。
\end{derivation}

偶子代数还把维数递推与 Spin 群直接联系起来。若 $p\ge1$，则有自然同构
\begin{equation}
\Cl_{p,q}^{0}
\cong
\Cl_{q,p-1}.
\label{eq:clifford-even-subalgebra}
\end{equation}
选取一个平方为 $+1$ 的单位向量 $e$。对 $e^\perp$ 中的向量 $v$，令
\[
\varphi(v)=ve.
\]
由于 $ve=-ev$，有
\[
(ve)^2=-v^2e^2=-Q(v),
\]
所以正负号交换；而所有 $ve$ 都属于偶子代数。再次由 泛性质 与维数比较便得到\eqnrefs{eq:clifford-even-subalgebra}。这解释了为什么 Spin 群虽然定义在 $\Cl_{p,q}$ 中，却完全生活在偶子代数里；也解释了偶数维手征旋量为何首先是 偶 Clifford 代数 的不可约表示。

从表示论角度，矩阵分类还有一个直接推论。若
\[
A\cong M_N(\mathbb D),
\qquad
\mathbb D\in\{\RR,\CC,\mathbb H\},
\]
则 简单 left module 可取为 $\mathbb D^N$，其交换子代数由 Schur 型结论给出
\[
\operatorname{End}_{A}(S)\cong\mathbb D^{\mathrm{op}}.
\]
当代数分裂成两个 简单 块 时，则出现两个不等价不可约模。因而“有几个不可约旋量模”“旋量维数是多少”“是否存在额外实或四元数结构”都可以从\eqnrefs{eq:real-clifford-classification} 系统读取，而无需逐个 符号差 猜 gamma 矩阵。

\section{Pin/Spin 群与正交群的二重覆盖}
\label{sec:pin-spin}
设 $B$ 非退化，$v\in V$ 且 $Q(v)\neq0$。在 Clifford 代数中
\[
v^{-1}=\frac{v}{Q(v)}.
\]
对 $x\in V$，定义 扭曲 共轭
\begin{equation*}\rho_v(x)=-vxv^{-1}.\end{equation*}
利用\eqnrefs{eq:clifford-anticommutator} 可得
\begin{equation}
\rho_v(x)
=x-2\frac{B(v,x)}{Q(v)}v,
\label{eq:orthogonal-reflection}
\end{equation}
这正是关于超平面 $v^\perp$ 的正交反射。Cartan--Dieudonn\'e 定理说明有限维非退化二次空间中的任意正交变换都可写成有限次此类反射的乘积，因此 Clifford 可逆元足以生成整个正交群。

\begin{derivation}[Clifford 共轭作用严格推出正交反射]
由 Clifford 基本关系
\[
vx+xv=2B(v,x)1
\]
得到
\[
vx=-xv+2B(v,x)1.
\]
右乘 $v^{-1}=v/Q(v)$（此处用到 $v$ 非迷向即 $Q(v)\neq0$，故 $v$ 可逆且 $v^{-1}=v/Q(v)$ 由 $v\cdot v/Q(v)=Q(v)/Q(v)=1$ 验证），有
\[
vxv^{-1}
=(-xv+2B(v,x)1)v^{-1}
=-xvv^{-1}+2B(v,x)v^{-1}
=-x+2\frac{B(v,x)}{Q(v)}v.
\]
第一项用到 $vv^{-1}=1$，第二项用到标量与代数元可交换。乘上 扭曲 共轭 定义中的负号，便得到\eqnrefs{eq:orthogonal-reflection}：
\[
\rho_v(x)=-vxv^{-1}=x-2\frac{B(v,x)}{Q(v)}v.
\]

再验证它保持二次型。写
\[
a=\frac{B(v,x)}{Q(v)},
\qquad
\rho_v(x)=x-2av.
\]
则由 $Q$ 与 $B$ 的关系 $Q(x)=B(x,x)$ 及 $B$ 双线性，
\[
\begin{aligned}
Q(\rho_v(x))
&=Q(x-2av)\\
&=B(x-2av,x-2av)\\
&=B(x,x)-2aB(x,v)-2aB(v,x)+4a^2B(v,v)\\
&=Q(x)-4aB(x,v)+4a^2Q(v)\\
&=Q(x)
-4\frac{B(v,x)^2}{Q(v)}
+4\frac{B(v,x)^2}{Q(v)}\\
&=Q(x).
\end{aligned}
\]
中间用到 $B(x,v)=B(v,x)$（对称双线性型）。所以 $\rho_v\in O(V,Q)$。

进一步验证 $\rho_v$ 确为反射：被反射的超平面为 $v^\perp=\{x:B(v,x)=0\}$。若 $x\in v^\perp$，则 $a=0$，$\rho_v(x)=x$，即超平面上的点不动；若 $x=v$，则 $a=B(v,v)/Q(v)=1$，$\rho_v(v)=v-2v=-v$，即 $v$ 方向反向。这正是关于 $v^\perp$ 的正交反射的几何定义。
\end{derivation}

例：$\Cl_2(\RR)$ 中取 $v=e_1$，$\rho_{e_1}(x)=x-2x_1e_1$ 即关于 $e_2$ 轴的反射；$\Cl_{1,3}(\RR)$ 中取 $v=\gamma^0$，$\rho_{\gamma^0}$ 反号空间分量而保持时间分量，正是宇称的代数实现，Pin 群与 Spin 群由此分别通过 扭曲 共轭 实现 $O(1,3)$ 与其双覆盖 $SO(1,3)$。公式要求 $v$ 可逆（$Q(v)\ne0$）：退化二次型上 $v^{-1}$ 不存在，需限制到非迷向锥外，复 Clifford 代数则所有非零向量可逆，公式全局成立。Cartan--Dieudonn\'e 定理进一步指出任意正交变换可表为不超过 $n$ 个反射之积，故 Clifford 可逆元经 扭曲 共轭 生成整个 $O(V,Q)$。
由满足 $Q(v)=\pm1$ 的单位向量在 Clifford 乘法下生成的群称为 Pin 群，记为 $\operatorname{Pin}(V,Q)$。其偶数个单位向量乘积构成 Spin 群
\[
\Spin(V,Q)
=\operatorname{Pin}(V,Q)\cap\Cl^0(V,Q).
\]
由 扭曲 共轭 得到群同态
\[
\rho:\operatorname{Pin}(V,Q)\to O(V,Q).
\]
限制到 偶 部分 后得到
\begin{equation}
1\longrightarrow\{\pm1\}
\longrightarrow\Spin(V,Q)
\overset{\rho}{\longrightarrow}SO(V,Q)
\longrightarrow1,
\label{eq:spin-double-cover}
\end{equation}
在合适的连通分支上它是二重覆盖。旋量不是“普通向量再加几个分量”，而是 Spin 群不可约表示空间中的元素；正交向量表示与旋量表示属于不同表示。

\begin{derivation}[为什么 Spin 映射的核是 $\{\pm1\}$]
分三步证明 $\ker\rho=\{\pm1\}$。

平凡作用推出中心化。设 $a\in\Spin(V,Q)$ 在向量表示上作用平凡，即
\[
 axa^{-1}=x,\qquad \forall x\in V.
\]
等价地，$ax=xa$ 对所有 $x\in V$ 成立，即 $a$ 与所有向量对易。由于 Clifford 代数 $\Cl(V,Q)$ 由 $V$ 作为代数生成（乘积和线性组合反复构造），$a$ 与整个 $\Cl(V,Q)$ 对易。因此 $a$ 属于 Clifford 代数的中心
\[
 Z(\Cl(V,Q))=\{c\in\Cl(V,Q):cb=bc,\;\forall b\in\Cl(V,Q)\}.
\]
严格地说：若 $a$ 与所有 $x\in V$ 对易，则与任意单项式 $x_1x_2\cdots x_k$ 对易（归纳法：$a(x_1\cdots x_k)=x_1(a x_2\cdots x_k)=\cdots=(x_1\cdots x_k)a$），因而与任意线性组合对易，即与整个代数对易。

中心的显式结构。对非退化二次型 $Q$，Clifford 代数的中心由维数奇偶性决定：
\begin{itemize}
 \item $\dim V$ 为偶数时，$Z(\Cl(V,Q))=\R\cdot 1$（仅标量）；
 \item $\dim V$ 为奇数时，$Z(\Cl(V,Q))=\R\cdot 1\oplus\R\cdot\omega$（标量加体积元 $\omega=e_1e_2\cdots e_n$）。
\end{itemize}
但 $\Spin(V,Q)\subset\Cl^{\rm even}(V,Q)$，而奇维情形的体积元 $\omega$ 属于奇子代数，故 $a$ 作为偶元素只能落在 $\R\cdot 1$ 中：$a=\lambda\cdot 1$，$\lambda\in\R$。

Clifford 范数 限定 $\lambda=\pm1$。Spin 群元素满足 Clifford 范数 $N(a)=a\tilde a=\pm1$（其中 $\tilde{}$ 是主对合，把每个生成元 $e_i\mapsto -e_i$）。对 $a=\lambda\cdot 1$，$\tilde a=\lambda\cdot 1$（标量被主对合保持），故
\[
 N(a)=\lambda^2=1,
\]
因此 $\lambda=\pm 1$。反过来，$(\pm1)x(\pm1)^{-1}=x$ 对任意 $x$ 成立，故 $\pm1\in\ker\rho$。因此
\[
 \ker\rho=\{\pm1\}.
\]
结合 Cartan--Dieudonn\'e 定理给出的满射性（每个正交变换是反射之积，而每个反射由 Spin 中元素通过 $\rho$ 实现），便得到\eqnrefs{eq:spin-double-cover}。
\end{derivation}

例：$\Spin(2)\cong U(1)$ 经 $\rho(e^{i\theta/2})$ 给出 $SO(2)$ 中转 $\theta$ 的旋转，核为 $\pm1$——即转 $2\pi$ 在 $\Spin(2)$ 中回到 $-1$，须转 $4\pi$ 才回到恒等；$\Spin(3)\cong SU(2)\cong S^3$ 经 $q\vec xq^{-1}$ 作用于纯虚四元数，核为 $q=\pm1$，是自旋 $1/2$ 粒子转 $2\pi$ 出负号的代数根源；$\Spin(4)\cong SU(2)\times SU(2)$，核为 $\{(1,1),(-1,-1)\}$，故 $SO(4)\cong(SU(2)\times SU(2))/\mathbb Z_2$。核 $\{\pm1\}$ 完全由 Clifford 代数的中心与 范数 决定，与矩阵实现无关：$\Spin(n)$（$n\ge3$）单连通，故是 $SO(n)$ 的万有覆盖，所有 $SO(n)$ 的投影表示提升为 $\Spin(n)$ 的真实表示。物理上，$\Spin$ 在旋量表示中于 $-1$ 处取负号，正是 Fermi 统计的几何来源；Spin 结构的存在性 $w_2(TM)=0$ 是 Dirac 算子定义与 Atiyah--Singer 指标公式的前提。
Spin 群的无穷小结构不是一个额外定义，而是 Clifford 偶子代数中的 双矢 部分。把 $\mathfrak{so}(V,B)$ 视为保持 $B$ 的反对称无穷小变换。对 $u,v\in V$ 定义
\begin{equation*}\Sigma(u\wedge v)
=\frac14(uv-vu)\in\Cl^0(V,Q).\end{equation*}
则它对向量的交换子作用为
\begin{equation}
[\Sigma(u\wedge v),w]
=B(v,w)u-B(u,w)v.
\label{eq:spin-bivector-so-action}
\end{equation}
右边正是 $u\wedge v$ 对应的标准 $\mathfrak{so}(V,B)$ 生成元。因此
\begin{equation}
\mathfrak{spin}(V,B)
\cong\Lambda^2V
\cong\mathfrak{so}(V,B)
\label{eq:spin-lie-algebra-isomorphism}
\end{equation}
作为 Lie 代数 同构；Spin 与 $SO$ 的差异是全局拓扑二重覆盖，而不是局部 Lie 代数 的差异。

\begin{derivation}[双矢 的 Clifford 交换子为何给出正交 Lie 代数，即\eqnrefs{eq:spin-bivector-so-action}]
利用 Clifford 关系 $vw=-wv+2B(v,w)$ 逐步展开 $[uv,w]$，证明 双矢 $u\wedge v$ 的 Clifford 交换子恰好是 $\mathfrak{so}(V,B)$ 的标准生成元。

计算 $[uv,w]$。由 Clifford 关系 $vw=-wv+2B(v,w)$、$wu=-uw+2B(u,w)$，把 $uvw=u(vw)$ 与 $wuv=(wu)v$ 中的相邻向量对展开：
\begin{align*}
 [uv,w]
 &=uvw-wuv
 =u(vw)-(wu)v
 =u(-wv+2B(v,w))-(-uw+2B(u,w))v
 =(-uwv+2B(v,w)u)+(uwv-2B(u,w)v)
 &=2B(v,w)u-2B(u,w)v,
\end{align*}
其中 $-uwv$ 与 $+uwv$ 直接相消，标量因子 $2B(\cdot,\cdot)$ 与代数元可交换而提出。

反对称化。交换 $u,v$ 后同理有
\begin{equation*}
 [vu,w]
 =2B(u,w)v-2B(v,w)u.
\end{equation*}
两式相减再除以 $4$（即取 $u\wedge v=(uv-vu)/2$ 的交换子），
\begin{align*}
 \left[\frac{uv-vu}{2},w\right]
 &=\frac{1}{2}\bigl([uv,w]-[vu,w]\bigr)\\
 &=\frac{1}{2}\bigl(2B(v,w)u-2B(u,w)v
   -2B(u,w)v+2B(v,w)u\bigr)\\
 &=2B(v,w)u-2B(u,w)v,
\end{align*}
再除以 $2$ 得
\begin{equation*}
 \left[\frac{uv-vu}{4},w\right]
 =B(v,w)u-B(u,w)v,
\end{equation*}
即\eqnrefs{eq:spin-bivector-so-action}。

验证反对称性。定义 $A_{u,v}(w)=B(v,w)u-B(u,w)v$。对任意 $x,w$，
\begin{equation*}
 B(A_{u,v}x,w)+B(x,A_{u,v}w)
 =B\bigl(B(v,x)u-B(u,x)v,\,w\bigr)+B\bigl(x,\,B(v,w)u-B(u,w)v\bigr)=0,
\end{equation*}
展开后
\[
B(v,x)B(u,w)-B(u,x)B(v,w)+B(v,w)B(x,u)-B(u,w)B(x,v)=0,
\]
其中 $B(v,x)B(u,w)$ 项与 $B(u,w)B(x,v)$ 项配对抵消（用 $B$ 对称性 $B(x,v)=B(v,x)$），$B(u,x)B(v,w)$ 项与 $B(v,w)B(x,u)$ 项配对抵消，故 $A_{u,v}\in\mathfrak{so}(V,B)$（关于 $B$ 的反对称线性变换）。

维数匹配。$\Lambda^2V$ 的维数为 $\binom{n}{2}=n(n-1)/2$，$\mathfrak{so}(V,B)$ 的维数也是 $n(n-1)/2$（反对称矩阵）。非退化 $B$ 下，映射 $u\wedge v\mapsto A_{u,v}$ 是单射（核为零：若 $A_{u,v}=0$，取 $w$ 使 $B(v,w)\neq0$ 则 $u\propto v$，但 $u\wedge v=0$），两个空间维数相同，故为同构，即\eqnrefs{eq:spin-lie-algebra-isomorphism}。
\end{derivation}

例：$\mathfrak{spin}(3)$ 与角动量。取 $V=\RR^3$ Euclidean，双矢 基 $e_2e_3,e_3e_1,e_1e_2$ 对应 $\Sigma_1=\frac{e_2e_3}{2},\ \Sigma_2=\frac{e_3e_1}{2},\ \Sigma_3=\frac{e_1e_2}{2}$，$[\Sigma_i,\Sigma_j]=\varepsilon_{ijk}\Sigma_k$ 正是 $\mathfrak{so}(3)$ 标准基的对易关系，乘以 $-i$ 后即角动量对易 $[J_i,J_j]=i\varepsilon_{ijk}J_k$。旋量表示下，四维 Minkowski 空间的双矢 基 $\frac14\gamma^\mu\gamma^\nu$（$\mu<\nu$）通过换位子给出 Lorentz 生成元 $\Sigma^{\mu\nu}$：空间旋转 $\Sigma^i=\frac12\varepsilon^{ijk}\Sigma^{jk}$ 与推进 $K^i=\Sigma^{0i}$ 满足 $\mathfrak{so}(1,3)$ 标准对易关系；Dirac 旋量经 $S=\exp(-\frac{i}{2}\omega_{\mu\nu}\Sigma^{\mu\nu})$ 变换，半角系数正来自 双矢--向量同构中的 $1/2$ 因子。整体上，每点切空间的 双矢 拼成 $\Lambda^2TM\cong\mathfrak{so}(T_pM)$，给出标架丛的 Spin 结构。

这一同构也解释旋转的半角。设 Euclidean 二维平面由正交单位向量 $e_1,e_2$ 张成，并令
\[
J=e_1e_2,
\qquad J^2=-1.
\]
取 rotor
\begin{equation*}R(\theta)=\exp\!\left(-\frac{\theta}{2}J\right)
=\cos\frac\theta2-J\sin\frac\theta2.\end{equation*}
向量通过 共轭 变换
\[
x\longmapsto R(\theta)xR(\theta)^{-1}.
\]
指数中的参数是 $\theta/2$，但向量空间中得到的实际旋转角是 $\theta$。于是 $R(2\pi)=-1$ 而 $R(4\pi)=1$：这不是三维 Pauli 矩阵的偶然性质，而是二重覆盖在 Lie 群 层面的直接表现。

取 $V=\RR^3$ 与 Euclidean 二次型。其 偶 Clifford 代数 与四元数同构：
\[
\Cl^0_3(\RR)\cong\mathbb H.
\]
单位四元数恰构成 $S^3$，并与 $SU(2)$ 同构。因而
\[
\Spin(3)
\cong SU(2),
\qquad
SO(3)\cong SU(2)/\{\pm I\}.
\]
上一节四元数半角公式因此不是独立技巧，而是一般 Spin 二重覆盖在三维 Euclidean 空间中的具体实现。

\section{旋量模与 gamma 矩阵作为 Clifford 表示}
\label{sec:spinor-gamma}
给定 Clifford 代数 $\Cl(V,Q)$，旋量空间 $S$ 是其某个不可约左模，亦即存在代数表示
\[
\pi:\Cl(V,Q)\to\operatorname{End}(S).
\]
向量 $v\in V$ 的像
\[
\gamma(v)=\pi(v)
\]
满足
\begin{equation*}\gamma(u)\gamma(v)+\gamma(v)\gamma(u)
=2B(u,v)I_S.\end{equation*}
若选基 $e_\mu$，记 $\gamma_\mu=\gamma(e_\mu)$，便得到熟悉的 gamma 矩阵关系。矩阵维数、实条件和 手征性 都由 $\Cl(V,Q)$ 的表示理论决定，而不是凭经验猜测。

偶数维 $n=2m$ 时可定义体积元对应的 手征性 算符。对适当归一化的正交基，令
\[
\Gamma_*=c_m\gamma_1\gamma_2\cdots\gamma_{2m},
\qquad
\Gamma_*^2=I.
\]
它与每个 $\gamma_a$ 反对易，并把旋量空间分解为
\[
S=S^+\oplus S^-,
\qquad
\Gamma_*\psi_\pm=\pm\psi_\pm.
\]
这就是 Weyl 分解的纯代数来源。
偶数维复二次空间的旋量模可以不依赖任何具体 gamma 矩阵显式构造。设
\[
V_{\CC}=W\oplus W^*
\]
是一个 极大 isotropic 极化，使 $W$ 与 $W^*$ 各自在 $B$ 下全等距为零，而二者通过自然配对耦合。定义
\begin{equation*}S=\Lambda^\bullet W^*.\end{equation*}
对 $\xi\in W^*$ 用 exterior multiplication 作用，对 $w\in W$ 用 收缩 作用：
\[
\gamma(\xi)=\sqrt2\,\varepsilon(\xi),
\qquad
\gamma(w)=\sqrt2\,\iota_w.
\]
因为 产生 与 湮灭 算符 满足
\[
\{\varepsilon(\xi),\varepsilon(\eta)\}=0,
\qquad
\{\iota_w,\iota_z\}=0,
\]
\[
\{\iota_w,\varepsilon(\xi)\}=\xi(w)I,
\]
它们正给出 Clifford 关系。这就是旋量空间的 费米子 Fock 模型：旋量不是“几何向量的复数版本”，而是 极大 isotropic 子空间 上的外代数。

若 $\dim_{\CC}V=2m$，则 $\dim W=m$，从而
\[
\dim_{\CC}S=2^m.
\]
外代数的奇偶分次进一步给出
\begin{equation}
S^+=\Lambda^{\mathrm{even}}W^*,
\qquad
S^-=\Lambda^{\mathrm{odd}}W^*,
\label{eq:spinor-fock-chirality}
\end{equation}
且
\[
\dim S^+=\dim S^-=2^{m-1}.
\]
在偶数维，Clifford 的每个向量作用都会改变 宇称，因此
\[
\gamma(v):S^\pm\to S^\mp.
\]
这正是 手征性 算符 与 gamma 矩阵反对易的结构原因。

\begin{derivation}[从 极大 isotropic 极化 构造不可约旋量模]
选取对偶基
\[
w_1,\ldots,w_m\in W,
\qquad
\xi^1,\ldots,\xi^m\in W^*,
\]
使 $\xi^a(w_b)=\delta^a_b$。在 $S=\Lambda^\bullet W^*$ 上令
\[
a_a=\iota_{w_a},
\qquad
a_a^\dagger=\varepsilon(\xi^a).
\]
它们满足 正则反交换关系。具体验证：对任意 $\eta\in\Lambda^\bullet W^*$，
\[
\iota_{w_a}\varepsilon(\xi^b)\eta+\varepsilon(\xi^b)\iota_{w_a}\eta
=\xi^b(w_a)\eta=\delta_{ab}\eta,
\]
此为外代数上收缩--外积的标准恒等式。其余两条
\[
\{a_a,a_b\}=0,
\qquad
\{a_a^\dagger,a_b^\dagger\}=0,
\qquad
\{a_a,a_b^\dagger\}=\delta_{ab}I.
\]
来自 $\wedge$ 与 $\iota$ 各自的反对易性。真空 $|0\rangle=1\in\Lambda^0W^*$ 被所有 $a_a$ 消去（$\iota_{w_a}1=0$），而
\[
a_{i_1}^\dagger\cdots a_{i_k}^\dagger|0\rangle
=\xi^{i_1}\wedge\cdots\wedge\xi^{i_k}
\]
生成整个 $S$。因此 Fock 基恰有
\[
\sum_{k=0}^m\binom mk=2^m
\]
个元素（二项式定理）。
\end{derivation}

例：$m=1$ 即二维复 Clifford，$S=\Lambda^\bullet W^*=\operatorname{span}\{1,\xi^1\}$，$a_1=\iota_{w_1}$、$a_1^\dagger=\varepsilon(\xi^1)$ 在基 $\{1,\xi^1\}$ 下为 $2\times2$ 幂零矩阵，$\gamma(w_1)=\sqrt2a_1$、$\gamma(\xi^1)=\sqrt2a_1^\dagger$ 满足二维 Clifford 关系，手征分解 $S^\pm$ 各一维；$m=2$ 即四维复 Clifford，Fock 基按 宇称 分组为 $S^+=\operatorname{span}\{1,\xi^1\wedge\xi^2\}$ 与 $S^-=\operatorname{span}\{\xi^1,\xi^2\}$，$\gamma(w_a)$ 把 $S^\pm$ 映到 $S^\mp$，正是 Dirac 旋量分解为左右手 Weyl 旋量的代数来源。量子场论中自由 Dirac 场的 Fock 空间即由此构造：$a_a^\dagger$ 为产生算子、$a_a$ 为湮灭算子，多费米子态 $a_{i_1}^\dagger\cdots a_{i_k}^\dagger|0\rangle$ 由 Pauli 不相容原理（$\{a_a^\dagger,a_b^\dagger\}=0$）保证每模式最多一个粒子——旋量场的反对易不是额外公理，而是 Clifford 表示在 Fock 模型下的自然结果。

等价地，也可以把不可约旋量模实现为 Clifford 代数的 极小左理想
\[
S\cong\Cl(V,Q)f,
\]
其中 $f$ 为合适的 primitive 幂等。Fock 模型更适合计算，极小-left-ideal 语言则更接近纯代数分类；两者描述的是同一个不可约表示结构。


上述 Fock 构造还可以看出一个容易混淆的事实：$S^\pm$ 一般不是完整 Clifford 代数的模，而是偶子代数的模。因为任意向量 Clifford 乘法都会交换 宇称，
\[
\gamma(v):S^+\leftrightarrow S^-.
\]
因此在偶数维，完整不可约 Clifford 模通常是
\[
S=S^+\oplus S^-,
\]
而 Spin 群由于位于偶子代数中，分别保持 $S^+$ 与 $S^-$。这正是“Dirac 旋量 可以分成两个 Weyl 旋量”背后的表示论语句；它首先是 偶-subalgebra 的分解，而不是动力学方程的结果。

体积元把这个结构编码得更紧凑。对复偶数维 $n=2m$，取归一化体积元
\[
\omega_{\CC}=c_m e_1e_2\cdots e_{2m},
\qquad
\omega_{\CC}^2=1.
\]
它与每个向量反对易，却与所有偶元素交换。因此
\[
P_\pm=\frac12(1\pm\omega_{\CC})
\]
是偶子代数作用下的投影，并给出
\[
S^\pm=P_\pm S.
\]
反过来，如果只保留 $\Cl_{2m}^0(\CC)$ 的表示，就看不到把 $S^+$ 与 $S^-$ 互换的 奇 生成元；这也是构造 Weyl 模时必须明确“表示的是 Spin/偶 Clifford，还是完整 Clifford”的原因。

实旋量的情况不能只把复矩阵取实部。设 $S_\RR$ 是一个不可约实 Clifford 模并复化为
\[
S_\CC=S_\RR\otimes_\RR\CC.
\]
其复表示是否仍不可约、是否分裂成互为共轭的模、以及是否携带平方为 $-1$ 的交换子，都由\eqnrefs{eq:real-clifford-classification} 中的 $\RR/\CC/\mathbb H$ 类型控制。更抽象地，若存在与 Clifford 作用交换的反线性算子 $J$，则
\[
J^2=+1
\]
给出 real 结构，而
\[
J^2=-1
\]
给出 四元数 结构。这样的 $J$ 是否存在以及其平方符号不是基选择问题，而是实表示类型的不变量。

\begin{derivation}[手征性为什么只在偶维给出两个 Spin 不可约分量]
令 $n=2m$，体积元记为 $\omega=e_1\cdots e_n$。对任意基向量 $e_a$，把 $e_a$ 积过其余 $n-1$ 个生成元得到
\[
e_a\omega=(-1)^{n-1}\omega e_a=-\omega e_a,
\]
最后等号用到 $n=2m$ 偶故 $n-1$ 奇。所以 $\omega$ 与所有 奇 Clifford 元反交换。若 $x$ 是偶元素，则由两个 奇 因子组成的每一项都贡献两次负号，因此
\[
[x,\omega]=0,
\qquad x\in\Cl^0.
\]
归一化后令 $\omega^2=1$（取 $\omega_\CC=c_m\omega$，$c_m=i^{-m^2}$ 或相应相位因子使 $\omega_\CC^2=1$），则 $P_\pm=(1\pm\omega_\CC)/2$ 满足
\[
P_\pm^2=\tfrac14(1\pm2\omega_\CC+\omega_\CC^2)=\tfrac12(1\pm\omega_\CC)=P_\pm,
\]
\[
P_+P_-=\tfrac14(1-\omega_\CC^2)=0,
\qquad
P_++P_-=1.
\]
故 $P_\pm$ 是互补正交投影，且
\[
S=P_+S\oplus P_-S=S^+\oplus S^-.
\]
因为 $\Spin(p,q)\subset\Cl_{p,q}^0$，Spin 作用与 $P_\pm$ 交换，因而分别保持两部分；但任意 $v\in V$ 满足
\[
P_\pm\gamma(v)=\tfrac12(1\pm\omega_\CC)\gamma(v)
=\tfrac12(\gamma(v)\pm\omega_\CC\gamma(v))
=\tfrac12(\gamma(v)\mp\gamma(v)\omega_\CC)
=\gamma(v)P_\mp,
\]
中间用到 $\omega_\CC\gamma(v)=-\gamma(v)\omega_\CC$（体积元与向量反交换），所以完整 Clifford 作用把两部分互换。奇数维时 $n-1$ 偶，体积元与每个向量交换而不是反交换，因此它区分的是两个 简单 Clifford 块（中心非平凡），而不产生同样的 Weyl 宇称 交换结构。
\end{derivation}

例：$n=2$ 时 $\omega_\CC=i e_1e_2$ 可对角化为 $\operatorname{diag}(1,-1)$，$S^\pm$ 各一维，是 $\Spin(2)\cong U(1)$ 的两个不同权表示；$n=4$ 时 $\omega_\CC=i\gamma^0\gamma^1\gamma^2\gamma^3=\gamma_5$，$P_\pm=\frac12(1\pm\gamma_5)$，四维 Dirac 旋量 $S=\CC^4$ 分解为各二维的 $S^\pm$。完整 Clifford 作用把 $S^+$ 映到 $S^-$，故向量耦合 $\bar\psi\gamma^\mu\psi$ 保持手征，而质量项 $m\bar\psi\psi=m(\bar\psi_L\psi_R+\bar\psi_R\psi_L)$ 必须混合左右手——这正是弱相互作用 V--A 结构与 Majorana 质量项的代数根源。Weyl 分解的存在性随维数 mod 8 周期变化：复情形周期为 $2$，实情形周期为 $8$（不可约模类型随 $(p-q)\bmod8$ 在 $\RR,\CC,\mathbb H$ 间循环），正是实 Clifford 代数 Bott 周期性的表示论体现。
四维复 Clifford 代数维数为 $2^4=16$。在一个四维不可约复旋量表示中，$\operatorname{End}(S)\cong M_4(\CC)$ 也有复维数 16，因此 Clifford 基元可以组织成
\begin{equation}
I,
\qquad
\gamma^\mu,
\qquad
\sigma^{\mu\nu}=\frac{\ii}{2}[\gamma^\mu,\gamma^\nu],
\qquad
\Gamma_*\gamma^\mu,
\qquad
\Gamma_*.
\label{eq:dirac-matrix-basis}
\end{equation}
维数分别为
\[
1+4+6+4+1=16.
\]
这组对象的存在首先是 Clifford 分次与外代数维数计数的结果；它与任何特定 Dirac 动力学无关。

在标准迹配对
\[
(A,B)_{\operatorname{tr}}
=\operatorname{tr}(AB)
\]
下，不同 Clifford 级 的基元具有相应正交性。因而任意 $4\times4$ 复矩阵都可在\eqnrefs{eq:dirac-matrix-basis} 上唯一展开。Fierz 重排的代数本质正是“在这个完备矩阵基上插入恒等分解”。

\begin{derivation}[四维 Clifford 基的维数与线性完备性]
从 $n=4$ 的 Clifford 代数维数公\eqnrefs{eq:clifford-dimension} 出发，按 级 计数：
\[
\dim\Lambda^0V=\binom40=1,
\qquad
\dim\Lambda^1V=\binom41=4,
\]
\[
\dim\Lambda^2V=\binom42=6,
\qquad
\dim\Lambda^3V=\binom43=4,
\qquad
\dim\Lambda^4V=\binom44=1.
\]
因此总维数为
\[
1+4+6+4+1=16.
\]
在四维不可约复旋量模上，
\[
\dim_\CC S=4,
\qquad
\dim_\CC\operatorname{End}(S)=4^2=16.
\]
复 Clifford 代数同构于 $M_4(\CC)$（由\eqnrefs{eq:real-clifford-classification} 的复化版本给出），所以忠实不可约表示把上述 16 个独立 Clifford 元映到 $\operatorname{End}(S)$ 的 16 个独立矩阵。维数相同立即说明它们构成整个矩阵空间的一组基。由此任何 $M\in M_4(\CC)$ 都有唯一展开
\[
M=aI+b_\mu\gamma^\mu
+\frac12 c_{\mu\nu}\sigma^{\mu\nu}
+d_\mu\Gamma_*\gamma^\mu+e\Gamma_*.
\]
系数可通过与相应基元取迹逐项投影得到。具体地，由正交性
\[
\operatorname{tr}(\gamma^\mu\gamma^\nu)=4\eta^{\mu\nu},\qquad
\operatorname{tr}(\sigma^{\mu\nu}\sigma^{\rho\sigma})=4(\eta^{\mu\rho}\eta^{\nu\sigma}-\eta^{\mu\sigma}\eta^{\nu\rho}),
\]
等，系数为
\[
a=\tfrac14\operatorname{tr}M,\qquad
b_\mu=\tfrac14\operatorname{tr}(\gamma_\mu M),\qquad
e=\tfrac14\operatorname{tr}(\Gamma_*M),
\]
等等。这正是 Fierz 重排恒等式的代数基础：把任意 $4\times4$ 矩阵展开为完备基后，可系统地重排旋量双线性。
\end{derivation}

例：Fierz 恒等式。把 $|v\rangle\langle w|$ 视为 $4\times4$ 矩阵并按 Clifford 基展开，再代入第二个双线性，得到 $\bar u\gamma^\mu v\,\bar w\gamma_\mu z=-\sum_\Gamma\bar u\Gamma z\,\bar w\Gamma v$（$\Gamma$ 遍历 $16$ 个基元），这是计算弱衰变振幅与 Muonium 寿命的标准工具。Dirac 矩阵的完备性还保证任意 Lorentz 协变顶角可写为 $\Gamma^\mu=a\gamma^\mu+b_\nu\sigma^{\mu\nu}q_\nu+\cdots$；Gordon 恒等式正是 $\gamma^\mu$ 在不同运动学基下的重展。对一般偶维 $n=2m$，基元数 $2^{2m}$ 与 $\dim M_{2^m}(\CC)=2^{2m}$ 一致；实 Clifford 代数则需区分矩阵类型，基元数仍为 $2^n$，矩阵规模与实/复/四元数结构随 $(p-q)\bmod 8$ 周期变化（Bott 周期性）。
本章得到的对象有清晰的依赖顺序：二次型 $Q$ 先定义 Clifford 代数，Clifford 可逆元再产生 Pin/Spin 群，Spin 群的不可约模才定义旋量。若把这个顺序倒置，直接从 Pauli 或 Dirac 矩阵开始，就会把“一个具体表示”误认为“母代数本身”。

在微分几何中，把每个切空间的 Clifford 代数拼成 Clifford 丛，并选择 Spin 结构 后，才能定义 旋量丛 与 Dirac-型 算符；在相对论量子理论中，gamma 矩阵则是 Lorentz Clifford 代数 的一个表示。两种应用都建立在同一个代数母结构之上，但它们的 符号差、实结构与动力学解释不同，不能在本章提前混为一谈。

上一段用 Cayley--Dickson 乘法
\[
(a,b)(c,d)=(ac-\bar d\,b,\;da+b\bar c),
\qquad a,b,c,d\in\mathbb H
\]
定义了八元数。为了看清它与四元数究竟差在哪里，先把共轭和范数补出来。定义
\[
\overline{(a,b)}=(\bar a,-b).
\]
于是
\[
\begin{aligned}
(a,b)\overline{(a,b)}
&=(a,b)(\bar a,-b)\\
&=(a\bar a-\overline{-b}\,b,\;-ba+b\overline{\bar a})\\
&=(|a|^2+|b|^2,0).
\end{aligned}
\]
因此八元数的范数为
\begin{equation*}\boxed{
N(a,b)=|a|^2+|b|^2
}\end{equation*}
并且非零八元数仍有逆元
\[
x^{-1}=\frac{\bar x}{N(x)}.
\]
所以失去结合律并不等于失去除法结构；八元数仍然是实数上的赋范可除代数。

更强的是范数具有乘法性。设
\[
x=(a,b),\qquad y=(c,d),
\]
则
\[
xy=(u,v),
\qquad
u=ac-\bar d b,
\qquad
v=da+b\bar c.
\]
于是
\[
N(xy)=|u|^2+|v|^2.
\]
把四元数范数写成 $|q|^2=q\bar q$，并利用
\[
|pq|=|p|\,|q|,
\qquad
\overline{pq}=\bar q\,\bar p,
\]
展开第一项：
\[
\begin{aligned}
|u|^2
&=|ac|^2+|\bar d b|^2
-2\operatorname{Re}\!\left(ac\,\overline{\bar d b}\right)\\
&=|a|^2|c|^2+|d|^2|b|^2
-2\operatorname{Re}(ac\bar b d).
\end{aligned}
\]
同理
\[
\begin{aligned}
|v|^2
&=|da|^2+|b\bar c|^2
+2\operatorname{Re}\!\left(da\,\overline{b\bar c}\right)\\
&=|d|^2|a|^2+|b|^2|c|^2
+2\operatorname{Re}(da c\bar b).
\end{aligned}
\]
四元数乘积的实部满足循环不变性
\[
\operatorname{Re}(q_1q_2q_3q_4)
=\operatorname{Re}(q_2q_3q_4q_1),
\]
所以两个交叉项相消。最终得到
\begin{equation*}\boxed{
N(xy)=N(x)N(y).
}\end{equation*}
这就是八元数仍属于实 composition 代数 的关键性质。

但八元数已经不再结合。取基
\[
\begin{aligned}
e_1&=(\mathbf i,0),&
 e_2&=(\mathbf j,0),&
 e_3&=(\mathbf k,0),\\
e_4&=(0,1),&
 e_5&=(0,\mathbf i),&
 e_6&=(0,\mathbf j),&
 e_7&=(0,\mathbf k).
\end{aligned}
\]
由 Cayley--Dickson 乘法直接算得
\[
e_1e_2=e_3,
\qquad
e_3e_4=e_7,
\]
因此
\[
(e_1e_2)e_4=e_7.
\]
另一方面
\[
e_2e_4=e_6,
\qquad
e_1e_6=-e_7,
\]
所以
\begin{equation*}\boxed{
(e_1e_2)e_4=-e_1(e_2e_4).
}\end{equation*}
这给出了一个完全显式的非结合例子。以后如果没有括号，三个八元数的乘积就没有唯一含义，不能像矩阵或四元数那样任意重新结合。

八元数虽然不结合，却满足交替律
\[
(xx)y=x(xy),
\qquad
 y(xx)=(yx)x.
\]
等价地，结合子
\[
[x,y,z]=(xy)z-x(yz)
\]
对三个变量是交替的；只要其中两个变量相同，结合子就为零。由 Artin 定理，任意两个八元数生成的子代数仍然结合。正因为如此，包含两个固定元素的许多计算仍可以像四元数一样进行，而真正的非结合性要到三个彼此独立的方向同时出现时才暴露出来。

这也解释了单位八元数与单位四元数的差别。单位四元数在普通乘法下组成群
\[
S^3\cong SU(2),
\]
而单位八元数虽然组成七维球面 $S^7$，却因为缺少结合律而不能成为通常意义下的群；它构成的是一个 Moufang 回路。因而不能从“单位八元数是七维球面”直接推断出一个普通 Lie 群结构。

现在再说明八元数与 $SU(3)$ 的联系究竟在哪里。八元数的实线性自同构群定义为
\[
G_2=\operatorname{Aut}(\mathbb O),
\]
它由保持八元数乘法的可逆实线性变换组成。因为自同构自动保持单位元和范数，$G_2$ 也保持七维虚八元数空间
\[
\operatorname{Im}\mathbb O\cong\mathbb R^7
\]
中的内积结构。

固定一个单位虚八元数 $e$，考虑稳定子
\[
(G_2)_e
=\{g\in G_2:g(e)=e\}.
\]
在与 $e$ 正交的六维实空间
\[
V=e^\perp\cap\operatorname{Im}\mathbb O
\]
上，左乘 $e$ 定义
\[
Jv=ev.
\]
由于 $e^2=-1$，并且由交替律可在只含 $e,v$ 的子代数中重新结合，
\[
J^2v=e(ev)=(e^2)v=-v.
\]
所以 $J^2=-I$，即 $J$ 给 $V$ 一个复结构；六维实空间因此可看成三维复空间。

若 $g\in(G_2)_e$，则它保持乘法且固定 $e$，于是
\[
g(Jv)=g(ev)=g(e)g(v)=e\,g(v)=Jg(v).
\]
所以稳定子中的变换都是复线性的。另一方面它们保持八元数范数给出的 厄米 内积，并且还保持由八元数乘法诱导的复体积形式。于是这个稳定子不是整个 $U(3)$，而恰好缩小为
\begin{equation*}\boxed{(G_2)_e\cong SU(3).}\end{equation*}
从维数也能看出结构是相容的：$\dim G_2=14$，单位虚八元数构成 $S^6$，若 $G_2$ 在其上传递，则稳定子的维数为
\[
14-6=8=\dim SU(3).
\]
维数计数本身不是同构证明，但与上面的复结构和体积形式说明完全一致。

因此正确的逻辑链是
\[
\mathbb O
\longrightarrow
G_2=\operatorname{Aut}(\mathbb O)
\longrightarrow
\text{固定一个单位虚八元数的稳定子 }SU(3),
\]
而不是
\[
\text{“八元数有八个基元”}
\Longrightarrow
\text{“八个 Gell--Mann 矩阵就是八元数”}.
\]
后者只比较了数字 $8$，忽略了乘法、结合性以及群与代数之间的层级差别。

Dirac 方程只有在 Lorentz 变换下保持形式不变，才真正是相对论方程。设坐标变换为
\[
x'^\mu=\Lambda^\mu{}_{\nu}x^\nu,
\qquad
\Lambda^Tg\Lambda=g.
\]
链式法则给出
\[
\partial'_\mu=(\Lambda^{-1})^\nu{}_{\mu}\partial_\nu.
\]
若旋量按
\[
\psi'(x')=S(\Lambda)\psi(x)
\]
变换，则新坐标中的 Dirac 方程为
\[
(\ii\hbar\gamma^\mu\partial'_\mu-mc)\psi'(x')=0.
\]
代入上式并左乘 $S^{-1}$，得到
\[
\ii\hbar S^{-1}\gamma^\mu S(\Lambda^{-1})^\nu{}_{\mu}\partial_\nu\psi-mc\,\psi=0.
\]
为了与原方程
\[
(\ii\hbar\gamma^\nu\partial_\nu-mc)\psi=0
\]
完全一致，必须满足
\begin{equation}
\boxed{
S^{-1}(\Lambda)\gamma^\mu S(\Lambda)
=\Lambda^\mu{}_{\nu}\gamma^\nu
}.
\label{eq:spinor-lorentz-intertwiner}
\end{equation}
因此旋量变换并不是把四个分量当成普通四矢量变换；它由一个不同的矩阵表示 $S(\Lambda)$ 实现，并通过上式与 Lorentz 矢量表示相联系。

看无穷小变换最清楚。写
\[
\Lambda^\mu{}_{\nu}
=\delta^\mu{}_{\nu}+\omega^\mu{}_{\nu}+O(\omega^2),
\]
其中 Lorentz 条件要求
\[
\omega_{\mu\nu}=-\omega_{\nu\mu}.
\]
定义
\[
\sigma^{\mu\nu}=\frac{\ii}{2}[\gamma^\mu,\gamma^\nu].
\]
由 Clifford 关系可以直接算得
\begin{equation*}\boxed{
[\sigma^{\mu\nu},\gamma^\rho]
=2\ii\bigl(g^{\nu\rho}\gamma^\mu-g^{\mu\rho}\gamma^\nu\bigr)
}.\end{equation*}
于是取
\[
S(\Lambda)
=I-\frac{\ii}{4}\omega_{\mu\nu}\sigma^{\mu\nu}+O(\omega^2),
\]
便有
\[
\begin{aligned}
S^{-1}\gamma^\rho S
&=\gamma^\rho
-\frac{\ii}{4}\omega_{\mu\nu}[\gamma^\rho,\sigma^{\mu\nu}]+O(\omega^2)\\
&=\gamma^\rho+\omega^\rho{}_{\nu}\gamma^\nu+O(\omega^2),
\end{aligned}
\]
正好满足 \eqnrefs{eq:spinor-lorentz-intertwiner}。积分无穷小变换得到
\[
\boxed{
S(\Lambda)
=\exp\!\left[-\frac{\ii}{4}\omega_{\mu\nu}\sigma^{\mu\nu}\right]
}.
\]
这就是 Dirac 旋量的 Lorentz 表示。

纯空间转动时，生成元可以写成
\[
\Sigma^k=\frac12\varepsilon^{kij}\sigma^{ij}
=\begin{pmatrix}
\sigma_k&0\\0&\sigma_k
\end{pmatrix},
\]
因此转角 $\boldsymbol\theta$ 对应
\[
S(R)=\exp\!\left(-\frac{\ii}{2}\boldsymbol\theta\cdot\boldsymbol\Sigma\right).
\]
这里再次出现半角。这与前面 $SU(2)\to SO(3)$ 的二重覆盖完全同源：Dirac 旋量中的自旋自由度正是 $SU(2)$ 基本表示在相对论时空中的延伸。
取平面波 Ansatz
\[
\psi(x)=u(p)\ee^{-\ii p\cdot x/\hbar},
\qquad p^\mu=(E/c,\mathbf p).
\]
代入 Dirac 方程得到
\begin{equation}
(\slashed p-mc)u(p)=0,
\qquad
\slashed p\equiv\gamma^\mu p_\mu.
\label{eq:dirac-u-equation}
\end{equation}
左乘 $\slashed p+mc$ 得
\[
(\slashed p+mc)(\slashed p-mc)u
=(p^2-m^2c^2)u=0,
\]
所以非零解要求
\[
p^2=m^2c^2,
\qquad
E=\pm\sqrt{c^2\mathbf p^2+m^2c^4}.
\]
这说明 Dirac 方程的一阶结构并没有改变相对论色散关系；它只是把二阶 Klein--Gordon 算符因式分解成矩阵值的一阶算符。

在 Dirac 表示中写
\[
u(p)=\begin{pmatrix}\varphi\\\chi\end{pmatrix},
\]
其中 $\varphi,\chi$ 都是二分量旋量。方程 \eqnrefs{eq:dirac-u-equation} 等价于
\[
(E-mc^2)\varphi-c(\boldsymbol\sigma\cdot\mathbf p)\chi=0,
\]
\[
c(\boldsymbol\sigma\cdot\mathbf p)\varphi-(E+mc^2)\chi=0.
\]
对正能量 $E>0$，第二式给出
\[
\chi=\frac{c\,\boldsymbol\sigma\cdot\mathbf p}{E+mc^2}\varphi.
\]
选取归一化二旋量 $\xi_s$（$s=1,2$），并采用
\[
\bar u_s(p)u_{s'}(p)=2mc\,\delta_{ss'},
\]
可取
\begin{equation*}\boxed{
u_s(p)
=\sqrt{\frac{E+mc^2}{c}}
\begin{pmatrix}
\xi_s\\[3pt]
\dfrac{c\,\boldsymbol\sigma\cdot\mathbf p}{E+mc^2}\xi_s
\end{pmatrix}
}.\end{equation*}
验证归一化时用到
\[
(\boldsymbol\sigma\cdot\mathbf p)^2=\mathbf p^2I
\]
以及
\[
E^2-c^2\mathbf p^2=m^2c^4.
\]
确实，
\[
\begin{aligned}
\bar u u
&=\frac{E+mc^2}{c}\left[1-\frac{c^2\mathbf p^2}{(E+mc^2)^2}\right]\\
&=\frac{(E+mc^2)^2-c^2\mathbf p^2}{c(E+mc^2)}\\
&=2mc.
\end{aligned}
\]

负频率支可以写成
\[
\psi(x)=v(p)\ee^{+\ii p\cdot x/\hbar},
\qquad E=+\sqrt{c^2\mathbf p^2+m^2c^4},
\]
于是满足
\[
(\slashed p+mc)v(p)=0.
\]
它在相对论量子论中对应反粒子自由度；Dirac 方程的四个分量并不是四个独立粒子态，而是由两个自旋自由度与正负频率结构共同组成。
质量壳上 $p^2=m^2c^2$ 时，
\[
(\slashed p-mc)(\slashed p+mc)=0.
\]
因此
\[
\Lambda_+(p)=\frac{\slashed p+mc}{2mc},
\qquad
\Lambda_-(p)=\frac{-\slashed p+mc}{2mc}
\]
满足
\[
\Lambda_+^2=\Lambda_+,
\qquad
\Lambda_-^2=\Lambda_-,
\qquad
\Lambda_+\Lambda_-=0,
\qquad
\Lambda_++\Lambda_-=I,
\]
其中第二个投影的具体写法随把负能量支写成 $p^0<0$ 还是用正能量反粒子 $v(p)$ 的约定而略有变化。对正能量 $u_s$ 最常用的是
\[
\Lambda_+u_s=u_s.
\]

在归一化 $\bar u_su_{s'}=2mc\,\delta_{ss'}$ 下，自旋完备关系为
\begin{equation*}\boxed{
\sum_{s=1}^{2}u_s(p)\bar u_s(p)=\slashed p+mc
}.\end{equation*}
为什么右端只能是这个形式，可以不靠逐分量计算看出。左端是作用在正能量解空间上的秩二矩阵，且每一项都满足
\[
(\slashed p-mc)u_s=0,
\]
所以
\[
(\slashed p-mc)\sum_su_s\bar u_s=0.
\]
Lorentz 协变性又要求它只能由 $I$ 与 $\slashed p$ 线性组合而成。结合质量壳条件和迹归一化，就唯一固定为 $\slashed p+mc$。同理
\[
\boxed{
\sum_{s=1}^{2}v_s(p)\bar v_s(p)=\slashed p-mc
}
\]
在标准 $v$ 旋量归一化下成立。

这些完备关系的意义很重要：一旦要对未观测的自旋自由度求和，旋量分量的求和就可以被转化为 $\gamma$ 矩阵迹运算。这也是 Clifford 代数真正进入相对论计算的地方。
由
\[
(\gamma^5)^2=I,
\qquad
(\gamma^5)^\dagger=\gamma^5
\]
可定义两个正交投影
\[
P_L=\frac{1-\gamma^5}{2},
\qquad
P_R=\frac{1+\gamma^5}{2}.
\]
它们满足
\[
P_L^2=P_L,
\qquad
P_R^2=P_R,
\qquad
P_LP_R=0,
\qquad
P_L+P_R=I.
\]
于是任意 Dirac 旋量都可唯一分解为
\[
\psi=\psi_L+\psi_R,
\qquad
\psi_{L,R}=P_{L,R}\psi.
\]
因为 $\{\gamma^5,\gamma^\mu\}=0$，有
\[
\gamma^\mu P_L=P_R\gamma^\mu,
\qquad
\gamma^\mu P_R=P_L\gamma^\mu.
\]
把 Dirac 方程分别投影，得到
\[
\ii\hbar\gamma^\mu\partial_\mu\psi_L=mc\,\psi_R,
\]
\[
\ii\hbar\gamma^\mu\partial_\mu\psi_R=mc\,\psi_L.
\]
所以质量项把左右手征分量耦合起来；当 $m=0$ 时两式完全解耦：
\[
\ii\hbar\gamma^\mu\partial_\mu\psi_L=0,
\qquad
\ii\hbar\gamma^\mu\partial_\mu\psi_R=0.
\]
这说明 $\gamma^5$ 并不是为了凑足 16 个矩阵才引入的，它直接把无质量 Dirac 理论分解成两个互不耦合的 Weyl 旋量。

16 个 Dirac 基元只有在能够把任意 $4\times4$ 复矩阵重新投影到这组基上时，才真正体现出“基”的意义。设
\[
\mathcal B
=\{I,\gamma^\mu,\sigma^{\mu\nu},\gamma^5\gamma^\mu,\gamma^5\}.
\]
它们共有
\[
1+4+6+4+1=16
\]
个独立方向，而复 $4\times4$ 矩阵空间本身也是 16 维，所以只要证明线性无关就构成完备基。最方便的工具是矩阵迹。

先证明
\begin{equation}
\boxed{
\operatorname{tr}(\gamma^\mu\gamma^\nu)=4g^{\mu\nu}
}.
\label{eq:gamma-trace-two}
\end{equation}
由 Clifford 关系
\[
\gamma^\mu\gamma^\nu+\gamma^\nu\gamma^\mu
=2g^{\mu\nu}I
\]
取迹，并利用迹的循环性
\[
\operatorname{tr}(\gamma^\mu\gamma^\nu)
=\operatorname{tr}(\gamma^\nu\gamma^\mu),
\]
得到
\[
2\operatorname{tr}(\gamma^\mu\gamma^\nu)
=2g^{\mu\nu}\operatorname{tr}I
=8g^{\mu\nu}.
\]
于是 \eqnrefs{eq:gamma-trace-two} 成立。

任何奇数个 $\gamma$ 矩阵的迹都为零。以一个 $\gamma^\mu$ 为例，由
\[
\gamma^5\gamma^\mu\gamma^5=-\gamma^\mu
\]
以及 $(\gamma^5)^2=I$，有
\[
\begin{aligned}
\operatorname{tr}(\gamma^\mu)
&=\operatorname{tr}(\gamma^5\gamma^5\gamma^\mu)\\
&=\operatorname{tr}(\gamma^5\gamma^\mu\gamma^5)\\
&=-\operatorname{tr}(\gamma^\mu),
\end{aligned}
\]
故
\[
\operatorname{tr}\gamma^\mu=0.
\]
同理，对任意奇数 $2n+1$，把一个 $\gamma^5$ 从乘积左端反对易移动到右端会产生 $(-1)^{2n+1}=-1$，因此迹等于其负值，只能为零。

四个 $\gamma$ 的迹也可以完全由 Clifford 关系化简。记
\[
T^{\mu\nu\rho\sigma}
=\operatorname{tr}(\gamma^\mu\gamma^\nu\gamma^\rho\gamma^\sigma).
\]
先交换前两个矩阵：
\[
\gamma^\mu\gamma^\nu
=2g^{\mu\nu}I-\gamma^\nu\gamma^\mu.
\]
于是
\[
T^{\mu\nu\rho\sigma}
=8g^{\mu\nu}g^{\rho\sigma}
-\operatorname{tr}(\gamma^\nu\gamma^\mu\gamma^\rho\gamma^\sigma).
\]
再利用迹循环性把 $\gamma^\nu$ 移到末端，并依次用反对易关系把它穿过 $\gamma^\sigma,\gamma^\rho$，最终得到
\begin{equation}
\boxed{
\operatorname{tr}(\gamma^\mu\gamma^\nu\gamma^\rho\gamma^\sigma)
=4\bigl(
 g^{\mu\nu}g^{\rho\sigma}
-g^{\mu\rho}g^{\nu\sigma}
+g^{\mu\sigma}g^{\nu\rho}
\bigr)
}.
\label{eq:gamma-trace-four}
\end{equation}

带 $\gamma^5$ 的最重要迹是
\begin{equation*}\boxed{
\operatorname{tr}(\gamma^5\gamma^\mu\gamma^\nu\gamma^\rho\gamma^\sigma)
=-4i\varepsilon^{\mu\nu\rho\sigma}
}\end{equation*}
在本讲义采用 $g=\operatorname{diag}(1,-1,-1,-1)$、$\varepsilon^{0123}=+1$ 的约定下成立。这个公式的符号会随 $\gamma^5$ 与 $\varepsilon^{0123}$ 约定改变，因此全书必须固定约定后再使用。

这些迹公式让自旋求和从“逐分量乘四分量旋量”变成了 Clifford 代数运算。例如出现
\[
\sum_su_s(p)\bar u_s(p)=\slashed p+mc
\]
后，任何自旋求和的双线性量都会转化为形如
\[
\operatorname{tr}
[(\slashed p'+mc)\Gamma_1(\slashed p+mc)\Gamma_2]
\]
的矩阵迹，再用 \eqnrefs{eq:gamma-trace-two,eq:gamma-trace-four} 化成 Lorentz 标量。这正是 Dirac 代数在实际计算中的核心价值。
16 维 Clifford 基还把旋量双线性量按 Lorentz 变换性质分类。对 Dirac 旋量 $\psi$，常见的五类是
\[
\bar\psi\psi,
\qquad
\bar\psi\gamma^\mu\psi,
\qquad
\bar\psi\sigma^{\mu\nu}\psi,
\qquad
\bar\psi\gamma^\mu\gamma^5\psi,
\qquad
\bar\psi\gamma^5\psi.
\]
它们分别对应标量、四矢量、反对称二阶张量、轴矢量和赝标量。

以四矢量流为例。旋量变换为
\[
\psi' =S\psi,
\]
而由
\[
S^\dagger\gamma^0=\gamma^0S^{-1}
\]
可得
\[
\bar\psi'=\bar\psi S^{-1}.
\]
于是
\[
\begin{aligned}
\bar\psi'\gamma^\mu\psi'
&=\bar\psi S^{-1}\gamma^\mu S\psi\\
&=\Lambda^\mu{}_{\nu}\bar\psi\gamma^\nu\psi.
\end{aligned}
\]
所以
\[
j^\mu=\bar\psi\gamma^\mu\psi
\]
确实按 Lorentz 四矢量变换。这样前面由 Dirac 方程推出的连续性方程
\[
\partial_\mu j^\mu=0
\]
不仅在某个惯性系成立，而是一个 Lorentz 协变的守恒律。
